% === IV. SIMULATION ARCHITECTURE AND DATA PIPELINE ===========================
\section{Simulation Architecture and Data Pipeline}
\label{sec:platform}

The pipeline is exercised end to end inside a single architecture: a Formula
Student vehicle model in CARLA acting as the plant, a C\texttt{++} ROS~2
communication layer that drives it and republishes its state, and an
identification stack that consumes that state. CARLA~\cite{carla2017} was
selected because it exposes per-wheel slip angle, normal load and tire force
rather than chassis states alone, and because its synchronous mode with a fixed
physics step (\SI{0.033}{\second} here) stamps every sample on the simulation
clock; those two properties are what make the identifiability claims of
Section~\ref{sec:method} checkable rather than merely
arguable~\cite{kaur2021simsurvey}.

\subsection{Plant: Formula AI vehicle in CARLA}
\label{subsec:platform}

The plant is the driverless Formula Student--class vehicle model, instantiated
in CARLA (server 0.9.16, Unreal Engine 4.26) on the \texttt{silverstone} map: a
four-wheeled, front-steered, all-wheel-drive rigid-body vehicle whose lateral
forces are generated by the PhysX wheeled-vehicle tire
model~\cite{physxvehicle} under a configured friction coefficient and lateral
stiffness, not by a Magic Formula. The simulator therefore has no ``true''
Pacejka coefficients to recover; instead the lateral forces PhysX computes at
each wheel are taken as ground truth, and the identified Magic Formula is
validated against those forces.

\begin{table}[t]
\caption{Vehicle configuration and measured plant properties.}
\label{tab:vehicle}
\centering
\scriptsize
\setlength{\tabcolsep}{4pt}
\begin{tabular}{@{}lll@{}}
\toprule
Quantity  & Value & Unit\\
\midrule
Body mass (configured) & 240.0 & \si{\kilo\gram}\\
Static wheel load (total)      & 2644.5 & \si{\newton}\\
Total (kerb) mass $m$\textsuperscript{a} & 269.6 & \si{\kilo\gram}\\
Front/rear static split        & 51.2 / 48.8 & \si{\percent}\\
Wheelbase $L$                  & 1.5334 & \si{\meter}\\
$l_{f}$ / $l_{r}$              & 0.7381 / 0.7953 & \si{\meter}\\
Yaw inertia $I_{z}$\textsuperscript{c} & 51.1 & \si{\kilo\gram\meter\squared}\\
Max road-wheel angle           & 16.0 & \si{\degree}\\
Configured wheel friction\textsuperscript{d} & 1.0 & ---\\
Road-surface factor & 0.70 & ---\\
Effective friction\textsuperscript{e} & 0.70 & ---\\
Realised peak axle friction    & 0.71 & ---\\
\bottomrule
\end{tabular}

\vspace{2pt}
{\raggedright\footnotesize \textsuperscript{a}Total static wheel load divided
by $g$; this is the mass every equation in this paper uses, not the configured
body mass. 
\textsuperscript{c}Assumed; see Section~\ref{subsec:assumptions}.
\textsuperscript{d}A scenario parameter, not a vehicle property.
\textsuperscript{e}Configured wheel friction multiplied by the road surface's
own coefficient, as reported per wheel. It is the effective value, not the
configured one, that the identifier is scored against
(Section~\ref{subsec:frictionfloor}).\par}
\end{table}

\subsection{The plant's tire curve, in closed form}
\label{subsec:planttire}

The identification target is worth stating exactly, because it is knowable
here. PhysX computes the lateral force of each wheel with
\texttt{PxVehicle\-Compute\-Tire\-Force\-Default}~\cite{physxvehicle}; at zero
longitudinal slip and zero camber that expression collapses to
\begin{equation}
F_{y}(\alpha)=\mathrm{sgn}(\alpha)\,\mu F_{z}\,S_{1}(K),\qquad
K=\frac{C_{\mathrm{lat}}|\tan\alpha|}{\mu F_{z}},
\label{eq:physx}
\end{equation}
with $C_{\mathrm{lat}}=F_{z}^{\mathrm{rest}}\,k_{y}\,
S_{1}\!\big(3\bar F_{z}/k_{x}\big)$, $\bar F_{z}=F_{z}/F_{z}^{\mathrm{rest}}$
the normalised load, $(k_{y},k_{x})$ the wheel's lateral-stiffness parameters,
and $S_{1}(K)=\min(1,\,K-K^{2}/3+K^{3}/27)$ the engine's smoothing function.

Two properties of \eqref{eq:physx} shape everything downstream. It saturates at
$\mu F_{z}$, first reached at $K=3$, and it \emph{never falls off} past that
peak, so a Magic Formula --- whose $E$ and post-peak decay exist precisely to
describe a falling branch --- can match it up to the peak and no further. Its
scale is set entirely by $\mu$ and $C_{\mathrm{lat}}$, both of which the
simulator publishes, so the curve of Table~\ref{tab:ref} is read off the plant
rather than fitted to it.

\begin{table}[t]
\caption{The plant's lateral tire curve: the parameters of \eqref{eq:physx}, read from the simulator's own wheel configuration and telemetry.}
\label{tab:ref}
\centering
\begin{threeparttable}
\scriptsize
\setlength{\tabcolsep}{4pt}
\begin{tabular}{@{}lll@{}}
\toprule
Quantity & Value & Unit\\
\midrule
Effective friction $\mu$ & 0.70 & ---\\
Lateral stiffness $k_{y}$         & 17.0 & \si{\per\radian}\\
Max normalised load $k_{x}$       & 2.0 & ---\\
Initial slope $C_{\mathrm{lat}}$ at $\bar F_{z}=1$ & 14.875\,$F_{z}^{\mathrm{rest}}$ & \si{\newton\per\radian}\\
\quad front axle, $F_{z,f}=\SI{1373}{\newton}$ & 20.4 & \si{\kilo\newton\per\radian}\\
\quad rear axle, $F_{z,r}=\SI{1274}{\newton}$ & 19.0 & \si{\kilo\newton\per\radian}\\
Peak reached at, $\bar F_{z}=1$, $\mu=0.70$ & 8.0 & \si{\degree}\\
\bottomrule
\end{tabular}
\begin{tablenotes}[flushleft]\scriptsize
\item The initial slope is set by the \emph{rest} load and by the normalised load through $S_{1}$, not by $F_{z}$ alone, and falls monotonically with $\bar F_{z}$. Every axle slope quoted in this paper is taken at $\bar F_{z}=1$ and at the static axle loads listed here.
\item The peak angle follows from $K=3$ in \eqref{eq:physx} and therefore moves with $\mu$: \SI{8.0}{\degree} at the effective $\mu=\num{0.70}$ of the runs reported here, \SI{12.0}{\degree} at the $\mu=\num{1.05}$ of the prior-fitting run.
\end{tablenotes}
\end{threeparttable}
\end{table}

That the closed form is the right one was checked rather than assumed: on 1804
ground-truth wheel samples from a live lap
($P_{99}(|\alpha|)=\SI{3.66}{\degree}$), each predicted from its own measured
$\alpha$, $F_{z}$ and $\mu$, \eqref{eq:physx} returns $R^{2}=\num{0.99}$ at
\SI{16.1}{\newton} front and \SI{14.7}{\newton} rear RMSE.
Section~\ref{subsec:resultstire} uses this curve to separate the error of the
identified fit from the error of representing \eqref{eq:physx} in the form
\eqref{eq:mf} that the controller requires.

\begin{figure*}[t]
\centering
\resizebox{0.86\textwidth}{!}{% =============================================================================
% Fig. 3 -- Carla_ASU_Bridge architecture.
%
% Vector TikZ, drawn at the paper's own font sizes. Included by
% sections/method.tex inside a figure* (full text width).
%
% Reading order, left to right: CARLA server -> bridge lifecycle node ->
% ROS 2 topic layer -> autonomy stack.
%   * orange  = the ground-truth path (feedback/tire_forces). It leaves the
%               simulator and reaches only the scoring/validation nodes; it is
%               deliberately NOT wired into the identification loop.
%   * blue    = mpc/update_params, the gated model interface. Drawn inside the
%               autonomy stack because it is a service between two of this
%               workspace's nodes, not part of the bridge's topic layer.
%
% Nodes are stacked with relative positioning (below=... of ...) so that a
% change of wording cannot make two boxes overlap.
% =============================================================================

\begin{tikzpicture}[
  font=\scriptsize,
  >={Stealth[length=1.7mm,width=1.5mm]},
  % --- block styles ---
  colhead/.style  = {font=\scriptsize\bfseries, align=center},
  blk/.style      = {draw, rounded corners=1pt, align=left, inner sep=2.4pt,
                     text width=#1, fill=white, line width=0.35pt},
  srv/.style      = {blk=26mm, fill=black!5},
  brg/.style      = {blk=33mm, fill=black!3},
  top/.style      = {blk=36mm, font=\scriptsize\ttfamily, fill=white},
  gt/.style       = {top, fill=orange!20, draw=orange!80!black, line width=0.7pt},
  cons/.style     = {blk=31mm, fill=black!3},
  note/.style     = {font=\tiny, align=left, inner sep=1.5pt, text width=36mm},
  % --- edge styles ---
  flow/.style     = {->, line width=0.5pt, black!65},
  gtflow/.style   = {->, line width=0.9pt, orange!80!black},
  gateflow/.style = {->, line width=0.9pt, blue!60!black},
  lbl/.style      = {font=\tiny, inner sep=1.2pt, fill=white, midway}
]

\def\gap{3.2mm}   % vertical gap between stacked boxes

% ============================ 1. CARLA server ============================
\node[srv, anchor=north west] (ego) at (0,0)
  {\textbf{Ego vehicle}\\[1pt]
   {\ttfamily vehicle.vehicle.}\\{\ttfamily asurt\_fsai}\\[1pt]
   PhysX wheeled-vehicle solver,\\per-wheel tire friction};
\node[srv, below=\gap of ego]  (sens)
  {\textbf{Sensor actors}\\[1pt] IMU, GNSS, cameras, LiDAR};
\node[srv, below=\gap of sens] (tick)
  {\textbf{World tick}\\[1pt] synchronous, $\Delta t=0.033$\,s};
\node[srv, below=\gap of tick] (wt)
  {\textbf{Wheel telemetry}\\[1pt]
   per-wheel slip angle,\\normal load, lateral force};

\begin{scope}[on background layer]
  \node[draw, dashed, rounded corners=2pt, inner sep=3.6pt, fill=black!2,
        fit=(ego)(sens)(tick)(wt)] (carla) {};
\end{scope}
\node[colhead, above=1.4mm of carla.north]
  {CARLA server\\{\scriptsize\mdseries UE 4.26 / PhysX}};

% ======================== 2. Carla_ASU_Bridge node ========================
\node[brg, anchor=north west] (ctrl) at (5.75,0)
  {\textbf{Control application}\\[1pt]
   forwards speed and steering to\\CARLA's \emph{native} Ackermann
   controller (server-side PID)};
\node[brg, below=\gap of ctrl] (sth)
  {\textbf{Per-sensor threads}\\[1pt]
   one publish thread each;\\optional dedicated clients};
\node[brg, below=\gap of sth] (tth)
  {\textbf{Telemetry thread}\\[1pt]
   10\,Hz on its own thread --- each\\tick costs blocking CARLA RPCs};
\node[brg, below=\gap of tth] (clk)
  {\textbf{\ttfamily sensor\_clock}\\[1pt]
   every stamp drawn from the\\simulation epoch, not wall time};

\begin{scope}[on background layer]
  \node[draw, dashed, rounded corners=2pt, inner sep=3.6pt, fill=black!1,
        fit=(ctrl)(sth)(tth)(clk)] (bridge) {};
\end{scope}
\node[colhead, above=1.4mm of bridge.north]
  {\textsc{Carla\_ASU\_Bridge}\\{\scriptsize\mdseries ROS 2 lifecycle node (C\texttt{++})}};

% ============================ 3. ROS 2 topics ============================
\node[top, anchor=north west] (tsub) at (11.1,0)
  {/drive\\control/tire\_friction\\control/drag\_coefficient};
\node[note, below=1.2mm of tsub] (nsub)
  {\emph{subscribed} --- commands and\\runtime physics overrides};

\node[top, below=2.6mm of nsub] (tstate)
  {/odom\\feedback/steering\_angle\\feedback/imu\\feedback/speed};
\node[note, below=1.2mm of tstate] (nstate)
  {\emph{published} --- the signals the\\identification actually consumes};

\node[gt, below=2.6mm of nstate] (tgt)
  {feedback/tire\_forces};
\node[note, below=1.2mm of tgt] (ngt)
  {\textbf{ground truth} --- per-wheel slip\\
   angle, normal load and lateral\\
   force, straight from PhysX; no\\
   analytic tire model in between};

\begin{scope}[on background layer]
  \node[draw, dashed, rounded corners=2pt, inner sep=3.6pt, fill=black!1,
        fit=(tsub)(nsub)(tstate)(nstate)(tgt)(ngt)] (topics) {};
\end{scope}
\node[colhead, above=1.4mm of topics.north]
  {ROS 2 topic layer\\{\scriptsize\mdseries all stamped from \texttt{/clock}}};

% =========================== 4. Autonomy stack ===========================
\node[cons, anchor=north west] (pp) at (16.6,0)
  {\textbf{\ttfamily pure\_pursuit}\\[1pt]
   model-free bootstrap; needs\\only the wheelbase};
\node[cons, below=\gap of pp] (sid)
  {\textbf{\ttfamily On-Track-SysID}\\[1pt]
   residual NN $\rightarrow$ virtual sweep\\$\rightarrow$ Pacejka fit};
\node[cons, below=11mm of sid] (mpc)
  {\textbf{\ttfamily mpc\_path\_tracking}\\[1pt]
   \texttt{acados} NMPC; holds the\\admissibility gates};
\node[cons, below=\gap of mpc] (mgr)
  {\textbf{\ttfamily adaptive\_}\\\textbf{\ttfamily controller\_manager}\\[1pt]
   PP\,$\rightarrow$\,MPC arbitration};
\node[cons, below=\gap of mgr] (bch)
  {\textbf{Benchmark nodes}\\[1pt]
   1/5/10-step prediction metrics,\\$F_y$ against ground truth};

\begin{scope}[on background layer]
  \node[draw, dashed, rounded corners=2pt, inner sep=3.6pt, fill=black!2,
        fit=(pp)(sid)(mpc)(mgr)(bch)] (stack) {};
\end{scope}
\node[colhead, above=1.4mm of stack.north]
  {Autonomy stack\\{\scriptsize\mdseries this workspace}};

% ================================ arrows =================================
% CARLA <-> bridge, over the C++ client library.
% The two column groups have different heights, so carla.east and bridge.west
% do NOT share a y. Pin both endpoints to one row (the "Sensor actors" row)
% with the |- / -| intersection syntax so the arrow is exactly horizontal.
\draw[flow, <->, line width=0.7pt]
      (carla.east |- sens) -- (bridge.west |- sens)
      node[lbl, above, pos=0.5, yshift=0.3mm]{\ttfamily libcarla\_client}
      node[lbl, below, pos=0.5, yshift=-0.3mm]{RPC + streams};

% CONVENTION: every inter-column arrow begins and ends on a dashed GROUP
% boundary, never on an inner box -- an arrow that pierces the dashed frame
% reads as if it bypassed the column. The row it belongs to is carried by the
% y coordinate (|- <box>), so the arrow still lines up with its own topic.

% bridge -> published topics
\draw[flow] (bridge.east |- tstate) -- (topics.west |- tstate)
      node[lbl, above, pos=0.5]{publish};
\draw[gtflow] (bridge.east |- tgt) -- (topics.west |- tgt);

% commands travel right to left, back into the bridge
\draw[flow] (topics.west |- tsub) -- (bridge.east |- tsub)
      node[lbl, above, pos=0.5]{apply};

% topics -> autonomy stack, and the /drive return path
\draw[flow] (topics.east |- tstate) -- (stack.west |- tstate)
      node[lbl, above, pos=0.5]{state feedback};
\draw[flow] (stack.west |- tsub) -- (topics.east |- tsub)
      node[lbl, above, pos=0.5]{\ttfamily /drive};

% ground truth reaches ONLY the scoring nodes: out of the topic column at the
% tire_forces row, down the gutter, into the stack at the benchmark row.
\coordinate (gtout)  at (topics.east |- tgt);
\coordinate (gtbend) at ($(gtout)+(0.6,0)$);
\draw[gtflow] (gtout) -- (gtbend) |- (stack.west |- bch);
\node[font=\tiny, orange!80!black, rotate=90, anchor=south, inner sep=1.4pt]
      at ($(gtbend)!0.5!(gtbend |- bch)$) {validation only};

% the gated model interface, inside the stack
\draw[gateflow] (sid.south) -- (mpc.north)
      node[lbl, pos=0.5, align=center]
      {\ttfamily mpc/update\_params\\\rmfamily\emph{gated}};

% ================================= legend ================================
% Sits under columns 1--2, which are the shorter pair.
\node[anchor=north west, font=\tiny, align=left, draw, rounded corners=1.5pt,
      inner sep=4pt, fill=white, text width=88mm]
      at ($(carla.south west)+(0,-0.5)$)
  {\tikz[baseline=-0.5ex]\draw[gtflow](0,0)--(0.5,0);\ \textbf{ground-truth path.}
   Used only to \emph{score} the identifier; never inside the identification loop.\\[1.5pt]
   \tikz[baseline=-0.5ex]\draw[gateflow](0,0)--(0.5,0);\ \textbf{gated interface.}
   An identified tire set is accepted only if its grip ceiling, understeer sign and
   critical speed clear the controller's requirements (Sec.~\ref{subsec:gates}).\\[1.5pt]
   \tikz[baseline=-0.5ex]\draw[flow](0,0)--(0.5,0);\ ROS 2 topic traffic.};
\end{tikzpicture}
}
\caption{Architecture of the \bridge{} ROS~2 interface. The bridge links
against CARLA's C\texttt{++} client, drives the ego vehicle through CARLA's
native Ackermann controller, and republishes state, sensors and per-wheel
telemetry on ROS~2 topics stamped from the simulation clock. The
\texttt{feedback/\allowbreak tire\_forces} path (highlighted) carries the
simulator's own per-wheel slip angle, normal load and lateral force, and is the
ground truth against which the identifier is scored.}
\label{fig:bridge}
\end{figure*}

\subsection{\textup{\bridge{}}: interfaces and telemetry}
\label{subsec:bridge}

\bridge{}~\cite{carlaasubridge} is a C\texttt{++} ROS~2 lifecycle node linked
directly against CARLA's client library (Fig.~\ref{fig:bridge}). It spawns and
configures the ego vehicle, drives its inputs over a single
\texttt{ack\-er\-mann\_msgs/AckermannDriveStamped} on \texttt{/drive}, and
republishes simulator state as standard ROS~2 topics; all identification traffic
crosses this interface. On the feedback path, one property is specific to this
architecture and the rest of this paper rests on it: the simulator's per-wheel
slip angle, normal load and lateral force are republished on the same
simulation clock as the odometry the identifier consumes, so an identified tire
model can be scored against the forces the plant actually generated rather than
a second model. Appendix~\ref{sec:iface} lists the full interface and the
body-frame and sign conventions at the simulator--ROS~2 boundary.

\subsection{Data acquisition and preprocessing}
\label{subsec:data}

A model-free pure-pursuit controller drives the vehicle around the reference
raceline; it requires no tire knowledge beyond the wheelbase and is therefore a
valid bootstrap, as in the baseline \cite{ontrack2025,forzaeth2024}. Odometry,
the achieved steering angle and the drive command are sampled into a ring buffer
by a \SI{50}{\hertz} timer for \SI{30}{\second}, gated on
$v_{x}>\SI{1}{\meter\per\second}$. The timer runs on the simulation clock, which
advances only at the server's \SI{0.033}{\second} step, so the achieved
acquisition rate is the step rate.

Two mismatches follow. The acquisition timer asks for a sample every
\SI{0.020}{\second}, \SI{65}{\percent} more often than the server produces one,
so a fraction of consecutive buffer entries are repeated states and the one-step
residual target carries a zero-order-hold artefact on those pairs. Separately,
the one-step model of Section~\ref{subsec:model} is discretised at
$T_{s}=\SI{0.035}{\second}$, \SI{6}{\percent} longer than the physics step.
$T_{s}$ is empirical and is carried because the identification proved
insensitive to it: at \SI{0.020}{\second} and \SI{0.035}{\second} the corrected
configuration returns the \emph{same} last accepted coefficient set to four
decimal places, and the inherited configuration keeps $D$ on its \num{2.0} bound
at both, so the box-edge signature is not an artefact of the mismatch. Results
are reported at $T_{s}=\SI{0.035}{\second}$. Preprocessing follows the baseline:
a zero-phase low-pass filter, then symmetry augmentation by mirroring
$(v_{y},\omega,\delta)\rightarrow(-v_{y},-\omega,-\delta)$ at unchanged $v_{x}$.
Re-identification retriggers every \SI{30}{\second} by default.

Only the deployment stage leaves this architecture: validation draws on
\texttt{feedback/\allowbreak tire\_forces} and \texttt{feedback/imu}, neither of
which enters the identification loop, so the estimator is scored against a
channel it did not observe, and Section~\ref{sec:results} can attribute a
closed-loop difference to the identified model rather than to the interface
carrying it.
